% =============================================================================
\chapter{Implémentation et Résultats}
\label{ch:implementation-resultats}
% =============================================================================

Ce chapitre présente une lecture \emph{implémentation-centrée} de FusionCore. Toutes les sections sont rédigées à partir des composants, scripts et interfaces effectivement présents dans le dépôt (daemon Rust, contrôleur Python, outillage de test, documentation d’exploitation). L’objectif est double : (i) expliciter précisément la mécanique du système ; (ii) proposer une méthode d’évaluation reproductible sans introduire de résultats non observés.

% =============================================================================
\section{Implémentation effective}
\label{sec:mise-en-oeuvre}
% =============================================================================

\subsection{Injection dans la chaîne Proxmox/QEMU}

L’intégration s’appuie sur un chemin explicite : Proxmox lance \texttt{kvm}, redirigé vers \texttt{omega-qemu-launcher}, qui prépare un backend mémoire \texttt{memfd}, injecte \texttt{omega-uffd-bridge.so} dans le QEMU réel, puis connecte le descripteur \texttt{userfaultfd} au processus agent via socket Unix. Cette architecture évite les modifications du noyau Linux et préserve la maintenabilité des mises à jour Proxmox~\cite{proxmox_doc,qemu_doc,kvm_doc}.

Deux conséquences techniques sont importantes :
\begin{itemize}
    \item le \textbf{chemin critique des défauts de page} reste en espace utilisateur, donc instrumentable par logs/métriques ;
    \item la \textbf{transparence côté VM} est conservée : la VM ne modifie ni son noyau ni ses applications pour exploiter le paging distant.
\end{itemize}

\subsection{Plan de données Rust : modules et responsabilités}

Le binaire \texttt{omega-daemon} regroupe les tâches locales suivantes :
\begin{enumerate}
    \item \textbf{Store réseau} (TCP/TLS, port 9100) : persistance et restitution de pages ;
    \item \textbf{EvictionEngine} : sélection de pages froides et envoi distant ;
    \item \textbf{FaultBus} : adaptation de l’agressivité d’éviction selon le rythme des défauts ;
    \item \textbf{QuotaRegistry} : contrôle des budgets mémoire VM avant écriture distante ;
    \item \textbf{API cluster} (9200) et \textbf{API contrôle} (9300) ;
    \item \textbf{ordonnanceurs locaux} : vCPU, I/O disque, budgets GPU, exécution migration.
\end{enumerate}

Le runtime asynchrone Tokio et les structures concurrentes shardées (notamment pour le store de pages) servent à limiter la contention sur l’état partagé lorsque les flux \texttt{PUT\_PAGE}/\texttt{GET\_PAGE} et la collecte métrique sont simultanés.

\subsection{Plan de contrôle Python : décision cluster}

Le contrôleur \texttt{omega-controller} opère une vue globale du cluster en interrogeant l’état exporté par les nœuds. Les décisions sont structurées en trois familles :
\begin{itemize}
    \item \textbf{admission} (faisabilité RAM/vCPU/GPU d’une nouvelle VM) ;
    \item \textbf{placement} (choix du nœud cible par score multi-critères) ;
    \item \textbf{migration} (live ou cold selon pression, activité VM et contraintes ressources).
\end{itemize}

Le composant \texttt{ResilientCollector} ajoute une logique de \emph{circuit-breaker} et de cache à durée limitée pour maintenir la continuité décisionnelle lors d’indisponibilités transitoires de nœuds.

\subsection{Flux fonctionnels principaux}

\paragraph{Eviction/rappel d’une page.}
\begin{enumerate}
    \item Détection de pression mémoire locale (\texttt{/proc/meminfo}, PSI).
    \item Sélection d’une page froide par l’algorithme CLOCK.
    \item Émission \texttt{PUT\_PAGE} vers un store pair (TLS).
    \item Libération locale de la page ; marquage « distante » côté tracker VM.
    \item À l’accès ultérieur : défaut \texttt{userfaultfd}, \texttt{GET\_PAGE}, puis \texttt{UFFDIO\_COPY}.
\end{enumerate}

\paragraph{Migration pilotée par pression.}
\begin{enumerate}
    \item Collecte d’état cluster (RAM, CPU, GPU, I/O, contraintes VM).
    \item Évaluation d’un nœud cible admissible.
    \item Déclenchement \texttt{qm migrate} (live/cold selon politique).
    \item Nettoyage post-migration : quotas, pages distantes, réservations locales.
\end{enumerate}

% =============================================================================
\section{Invariants techniques vérifiables}
\label{sec:invariants-verifiables}
% =============================================================================

L’analyse croisée code + documentation opérationnelle permet d’énoncer quatre invariants techniques :
\begin{itemize}
    \item \textbf{I1 — Borne mémoire par VM :} les écritures de pages distantes sont validées contre un quota explicite ;
    \item \textbf{I2 — Rôles symétriques des nœuds :} un même nœud peut simultanément émettre et héberger des pages ;
    \item \textbf{I3 — Séparation plan critique / plan décisionnel :} le transport de pages reste découplé des API de pilotage ;
    \item \textbf{I4 — Dégradation contrôlée :} en cas d’échec de collecte, le contrôleur bascule sur un état cache borné temporellement.
\end{itemize}

Ces invariants ne constituent pas une preuve formelle complète, mais des propriétés observables et testables via les scénarios du dépôt.

% =============================================================================
\section{Protocole d’évaluation reproductible}
\label{sec:protocole-evaluation}
% =============================================================================

\subsection{Principe méthodologique}

La démarche expérimentale suit le principe \og scripts-first \fg{} : toute mesure reportée doit être reproductible par un script, une commande ou un endpoint explicitement disponible dans le projet (\texttt{docs/experiments.md}, \texttt{docs/metrics.md}, scripts \texttt{omega-lab}, \texttt{run\_stress}, \texttt{collect\_baseline}).

\subsection{Scénarios retenus et variables observées}

\begin{table}[htbp]
    \centering
    \caption{Scénarios, instrumentation et critères d’interprétation.}
    \label{tab:scenarios-eval}
    \begin{tabular}{p{0.13\linewidth} p{0.27\linewidth} p{0.25\linewidth} p{0.25\linewidth}}
        \toprule
        \textbf{ID} & \textbf{Procédure} & \textbf{Instrumentation} & \textbf{Critère observé} \\
        \midrule
        S1 & \texttt{node-a-agent --mode demo} & logs agent + store & intégrité lecture/écriture pages \\
        S2 & stress mémoire local & \texttt{/proc/meminfo} + PSI & activation éviction et variation \texttt{MemAvailable} \\
        S3 & rappel de pages sous charge & logs \texttt{page fault}/\texttt{GET\_PAGE} & distribution de latence de rappel \\
        M1 & migration live déclenchée API & sortie \texttt{qm migrate --online} & succès opérationnel + continuité service \\
        M2 & migration cold (VM idle/pression) & état pre/post migration & nettoyage ressources source \\
        R1 & indisponibilité d’un store & logs collecteur + cache TTL & maintien du pilotage cluster \\
        \bottomrule
    \end{tabular}
\end{table}

\subsection{Familles de métriques}

\begin{itemize}
    \item \textbf{Store distant :} \texttt{pages\_stored}, \texttt{put\_count}, \texttt{get\_count}, \texttt{hit\_rate\_pct}, connexions actives.
    \item \textbf{Agent/daemon :} \texttt{fault\_count}, \texttt{fault\_served}, \texttt{fault\_errors}, \texttt{pages\_evicted}, \texttt{pages\_fetched}.
    \item \textbf{Système :} cgroups v2 (CPU/I/O), PSI mémoire/I/O, \texttt{/proc/meminfo}.
    \item \textbf{Export Prometheus :} endpoint \texttt{/control/metrics} pour consolidation temporelle.
\end{itemize}

% =============================================================================
\section{Analyse technique et discussion}
\label{sec:resultats}
% =============================================================================

\subsection{Comportement du paging distant}

Les scénarios de validation montrent une chaîne cohérente éviction\,$\rightarrow$\,stockage distant\,$\rightarrow$\,rappel à la demande. L’effet dominant sur la performance perçue est la latence réseau (aller-retour + contention), devant le coût cryptographique TLS dans le contexte GbE visé par le projet.

\subsection{Interaction entre ordonnancement local et migration}

FusionCore applique une stratégie hiérarchique : ajustement local d’abord (vCPU, \texttt{cpu.weight}, \texttt{io.weight}, budgets GPU), puis migration si la pression persiste. Cette séquence réduit le risque d’oscillation (migrations répétées) et préserve les migrations pour les cas où l’élasticité locale devient insuffisante.

\subsection{Robustesse en présence de pannes partielles}

Le circuit-breaker du collecteur évite qu’un nœud temporairement indisponible bloque le contrôleur. Le compromis associé est explicite : pendant la fenêtre de cache, la décision peut s’appuyer sur un état partiellement obsolète. Cette limite est acceptable en mode dégradé, mais doit être intégrée au réglage des seuils de migration.

% =============================================================================
\section{Limites actuelles et perspectives}
\label{sec:discussion-limites}
% =============================================================================

\begin{itemize}
    \item \textbf{Sensibilité réseau :} la qualité du paging distant dépend fortement de la latence et de la gigue inter-nœuds.
    \item \textbf{Coût d’observabilité :} l’analyse post-mortem requiert une corrélation multi-sources (agent, daemon, controller, système).
    \item \textbf{Paramétrage dépendant du contexte :} les seuils d’éviction/migration doivent être calibrés pour chaque cluster réel.
    \item \textbf{Hétérogénéité GPU :} la contrainte de topologie matérielle réduit parfois fortement l’espace de placement.
\end{itemize}

En résumé, l’implémentation actuelle valide la faisabilité d’une mutualisation multi-ressources en espace utilisateur dans Proxmox, avec un protocole d’évaluation reproductible. La consolidation scientifique complète (comparaison quantitative inter-architectures, étude statistique longue durée) constitue une étape ultérieure qui doit être conduite sur banc instrumenté stable.
